\section{Appendix A. Reproducible Search Log and Database-Specific Queries}
\label{appendix:search_log}

This appendix provides the exact, database-specific queries, execution dates, filters, and retrieval statistics used for the structured search described in Section~\ref{sec:eligibility}, to allow the search to be independently audited, reproduced, or extended to additional databases.

\subsection{PubMed/MEDLINE}

\textbf{Execution date:} 2026-05-25. \textbf{Interface:} NCBI Entrez E-utilities (\texttt{esearch}), via the Biopython \texttt{Bio.Entrez} client.

\textbf{Filters applied:} publication date (\texttt{pdat}) restricted to 2000--3000 (i.e., 2000 onward, with no practical upper bound other than the retrieval date), and language restricted to English (\texttt{english[lang]}). No study-type, article-type, or subject-heading filters were applied.

\begin{quote}
\small\texttt{(angiograph*[tiab] OR fluoroscop*[tiab] OR "digital subtraction angiography"[tiab] OR DSA[tiab]) AND ("deep learning"[tiab] OR "machine learning"[tiab] OR CNN[tiab] OR transformer[tiab] OR "artificial intelligence"[tiab] OR AI[tiab]) AND ("2000"[pdat] : "3000"[pdat]) AND english[lang]}
\end{quote}

\textbf{Retrieved:} 4{,}071 records before deduplication; 3{,}312 unique records after within-source and cross-source deduplication (759 duplicates removed).

\subsection{Semantic Scholar}

\textbf{Execution date:} 2026-05-25. \textbf{Interface:} Semantic Scholar Graph API v1 (\texttt{/paper/search} endpoint).

\textbf{Filters applied:} publication year restricted to 2000--2026. No venue, study-type, or field-of-study filters were applied. Because the Semantic Scholar API does not support the same field-tagged Boolean syntax as PubMed, the search was decomposed into 21 shorter, relaxed free-text queries combining an imaging term (angiography, DSA, or fluoroscopy) with an AI/ML term, executed independently and merged. This design choice trades query precision for recall, consistent with the goal of maximizing sensitivity during automated retrieval and deferring precision to the downstream LLM-assisted screening stage.

\begin{quote}
\small\texttt{
deep learning angiography; deep learning digital subtraction angiography; machine learning angiography; deep learning fluoroscopy; artificial intelligence angiography; CNN angiography; transformer angiography; deep learning DSA; neural network DSA; machine learning DSA; artificial intelligence DSA; machine learning fluoroscopy; neural network fluoroscopy; artificial intelligence fluoroscopy; CNN fluoroscopy; deep learning vessel segmentation angiography; deep learning vascular imaging; neural network vascular segmentation; deep learning cerebrovascular angiography; deep learning coronary angiography; deep learning peripheral angiography; deep learning retinal angiography
}
\end{quote}

\textbf{Retrieved:} 1{,}758 records before deduplication; 1{,}753 unique records after deduplication (5 duplicates removed).

\subsection{Deduplication and Aggregation}

Records from both sources were pooled and deduplicated first by DOI (case-insensitive exact match) and, for records lacking a DOI, by normalized title (lowercased, punctuation-stripped exact match). This yielded 5{,}064 unique records from an initial pool of 5{,}829 (765 duplicates removed: 759 within/across PubMed, 5 within/across Semantic Scholar, with the remainder resolved at aggregation), matching the counts reported in Section~\ref{sec:eligibility}.

\subsection{Sources Not Queried}

\textbf{arXiv} was not queried directly. This was a deliberate scope decision to prioritize peer-reviewed venues over unreviewed preprints, rather than an oversight; an arXiv query string was drafted and is retained in the fetcher codebase for future extensions of this search but was not executed for the present review. \textbf{Scopus} and \textbf{Web of Science} were also not queried as part of the automated retrieval pipeline; we note this explicitly as a scope limitation rather than a claim of comprehensive database coverage. As reported in Section~\ref{sec:eligibility}, a retrospective venue audit of the 206 records reaching full-text extraction found that Semantic Scholar's aggregation nonetheless surfaced substantial IEEE (19 records, 9.2\%) and SPIE (15 records, 7.3\%) content, but only a single MICCAI/LNCS record and no ACM-venue records, indicating that coverage of ACM- and MICCAI-specific venues remains a limitation of the present search strategy.

\section{Appendix B. Stage 1 Screening Prompts}
\label{appendix:screening_prompts}

This appendix provides the prompts used during the title--abstract screening stage of the systematic review. Screening was implemented using a two-gate architecture. The first gate evaluated whether a study satisfied all predefined inclusion criteria. Only studies passing the inclusion gate were evaluated using the exclusion gate. The model was instructed to base all decisions solely on information explicitly available in the title and abstract.

\subsection{Prompt A: Inclusion Screening}

\begin{quote}
You are a systematic review assistant performing the \textbf{inclusion screening stage}.

Your task is to determine whether a research article satisfies all inclusion criteria for a systematic review of computational methods applied to angiographic image sequences.

\medskip

\textbf{Review Scope}

The review focuses on computational methods, including deep learning, machine learning, and traditional image processing approaches, applied to temporally ordered Digital Subtraction Angiography (DSA) or fluoroscopic image sequences.

Representative tasks include:

\begin{itemize}
\item Image enhancement and denoising
\item Artifact reduction
\item Motion correction and registration
\item Vessel segmentation
\item Detection and localization
\item Classification
\item Temporal sequence analysis
\item Outcome prediction
\end{itemize}

\medskip

\textbf{Inclusion Criteria}

All of the following criteria must be satisfied:

\begin{description}
\item[I1:] The imaging modality is DSA or fluoroscopy and contains temporally ordered image sequences.
\item[I2:] The study addresses at least one computational image analysis task relevant to angiographic imaging.
\item[I3:] The methodology is described with sufficient technical detail to identify the principal computational approach.
\item[I4:] Quantitative evaluation results are reported.
\end{description}

\medskip

\textbf{Decision Rules}

\begin{itemize}
\item \textbf{Pass}: All four inclusion criteria are clearly satisfied.
\item \textbf{Fail}: One or more inclusion criteria are not satisfied.
\item \textbf{Uncertain}: The title and abstract do not provide sufficient information to confidently determine eligibility.
\end{itemize}

\medskip

Base your assessment solely on information explicitly stated in the title and abstract. Do not infer information that is not clearly described.

For each article, provide:

\begin{enumerate}
\item A screening decision (\texttt{Pass}, \texttt{Fail}, or \texttt{Uncertain})
\item The inclusion criteria satisfied
\item The inclusion criteria not satisfied
\item A concise rationale supporting the decision
\item An indication of whether human review is recommended
\end{enumerate}
\end{quote}

\subsection{Prompt B: Exclusion Screening}

\begin{quote}
You are a systematic review assistant performing the \textbf{exclusion screening stage}.

The article has already satisfied all inclusion criteria. Your task is to determine whether any exclusion criterion applies.

\medskip

\textbf{Exclusion Criteria}

A study is excluded if any of the following criteria are satisfied:

\begin{description}
\item[E1:] The study focuses exclusively on CT, MRI, ultrasound, or static radiographic imaging without DSA or fluoroscopy.
\item[E2:] The study analyzes only a single static image and does not incorporate temporally ordered image sequences.
\item[E3:] The article is a review, commentary, editorial, or clinical narrative without an original computational contribution.
\item[E4:] The publication is available only as an abstract, poster, or brief summary without sufficient methodological detail.
\item[E5:] The article is not available in English.
\end{description}

\medskip

\textbf{Decision Rules}

\begin{itemize}
\item \textbf{Pass}: No exclusion criteria are triggered.
\item \textbf{Fail}: One or more exclusion criteria are triggered.
\item \textbf{Uncertain}: Insufficient information is available to determine whether an exclusion criterion applies.
\end{itemize}

\medskip

Base your assessment solely on information explicitly stated in the title and abstract. Do not infer information that is not clearly described.

For each article, provide:

\begin{enumerate}
\item A screening decision (\texttt{Pass}, \texttt{Fail}, or \texttt{Uncertain})
\item The exclusion criteria triggered
\item A concise rationale supporting the decision
\item An indication of whether human review is recommended
\end{enumerate}
\end{quote}

\subsection{Screening Decision Logic}

The final screening outcome was determined according to the following decision framework:

\begin{itemize}
\item Inclusion Pass + Exclusion Pass $\rightarrow$ Include
\item Inclusion Fail $\rightarrow$ Exclude
\item Inclusion Pass + Exclusion Fail $\rightarrow$ Exclude
\item Any Uncertain Decision $\rightarrow$ Human Review
\end{itemize}

All screening was performed with deterministic decoding (\emph{temperature}=0) to maximize reproducibility and minimize decision variability across records.
